\documentclass[12pt,a4paper]{article}
\usepackage[utf8]{inputenc}
\usepackage[bosnian]{babel}
\usepackage{amsmath}
\usepackage{amsfonts}
\usepackage{amssymb}
\usepackage{graphicx}
\usepackage{listings}
\usepackage{xcolor}
\usepackage{geometry}
\usepackage{hyperref}
\usepackage{float}

\geometry{margin=2.5cm}

% Definisanje stila za Julia kod
\lstdefinestyle{julia}{
    language=Python,
    basicstyle=\ttfamily\small,
    keywordstyle=\color{blue}\bfseries,
    commentstyle=\color{gray}\itshape,
    stringstyle=\color{red},
    numbers=left,
    numberstyle=\tiny\color{gray},
    stepnumber=1,
    numbersep=5pt,
    backgroundcolor=\color{white},
    showspaces=false,
    showstringspaces=false,
    showtabs=false,
    frame=single,
    tabsize=4,
    captionpos=b,
    breaklines=true,
    breakatwhitespace=false,
    escapeinside={\%*}{*)},
    morekeywords={function,end,if,else,elseif,for,while,return,using,println,error}
}

\title{\textbf{Zadaća 7 - Minimalno povezujuće stablo} \\ 
\large Rješavanje pomoću linearnog programiranja}
\author{Bakir Činjarević 19705 \& Amar Handanagić 19089}
\date{}

\begin{document}

\maketitle

\section{Uvod}

Cilj ovog zadatka je implementacija funkcije \texttt{min\_tree(W)} koja pronalazi minimalno povezujuće stablo (MST - Minimum Spanning Tree) pomoću linearnog programiranja u programskom jeziku Julia, koristeći pakete JuMP i GLPK.

\section{Matematički model}

Problem minimalnog povezujućeg stabla se može formulirati kao problem linearnog programiranja na sljedeći način:

\subsection{Varijable odlučivanja}

Za svaku granu $(i,j)$ u grafu (gdje $i < j$), definišemo binarnu varijablu:
\begin{equation}
x_{ij} \in \{0, 1\}
\end{equation}

gdje $x_{ij} = 1$ ako je grana $(i,j)$ uključena u minimalno stablo, a $x_{ij} = 0$ inače.

\subsection{Funkcija cilja}

Cilj je minimizirati ukupnu težinu stabla:
\begin{equation}
\min \sum_{i=1}^{n} \sum_{j=i+1}^{n} w_{ij} \cdot x_{ij}
\end{equation}

gdje je $w_{ij}$ težina grane između čvorova $i$ i $j$.

\subsection{Ograničenja}

Problem ima dva tipa ograničenja:

\subsubsection{Ograničenje broja grana}

Stablo sa $n$ čvorova mora imati tačno $n-1$ grana:
\begin{equation}
\sum_{i=1}^{n} \sum_{j=i+1}^{n} x_{ij} = n - 1
\end{equation}

\subsubsection{Ograničenja povezanosti (eliminacija ciklusa)}

Za svaki neprazan podskup $S$ čvorova (osim punog skupa), broj grana između $S$ i komplementa $S$ mora biti najmanje 1. Ovo osigurava povezanost grafa i eliminiše cikluse:
\begin{equation}
\sum_{i \in S, j \notin S} x_{ij} \geq 1 \quad \forall S \subset \{1, 2, \ldots, n\}, S \neq \emptyset, S \neq \{1, 2, \ldots, n\}
\end{equation}

\section{Implementacija}

Funkcija \texttt{min\_tree(W)} je implementirana u programskom jeziku Julia koristeći pakete JuMP i GLPK. Ključni dijelovi implementacije su:

\subsection{Kreiranje modela i varijabli}

Varijable su definisane kao cjelobrojne (Int) kako bi se osiguralo da algoritam vraća cjelobrojno rješenje, što je važno jer algoritmi unutrašnje tačke mogu dati necjelobrojno optimalno rješenje čak i kada postoji cjelobrojno optimalno rješenje:

\begin{lstlisting}[style=julia]
@variable(model, x[i=1:n, j=1:n; i < j && W[i,j] < Inf], Int, lower_bound=0, upper_bound=1)
\end{lstlisting}

\subsection{Generisanje ograničenja povezanosti}

Ograničenja povezanosti se generišu za sve neprazne podskupove čvorova (osim punog skupa). Za graf sa $n$ čvorova, postoji $2^n - 2$ takvih podskupova. Za mali broj čvorova (npr. $n \leq 8$), ovo je izvodljivo.

\subsection{Optimizacija i ekstrakcija rješenja}

Nakon optimizacije, izvuku se grane za koje je $x_{ij} = 1$, sortiraju se po prvom pa po drugom čvoru za konzistentan prikaz, i formira se matrica $T$ gdje svaki red predstavlja jednu granu.

\section{Rezultati testiranja}

Funkcija je testirana na četiri različita primjera:

\subsection{Testni primjer 1 - Primjer iz zadatka}

Ulazna težinska matrica:
\begin{equation}
W = \begin{bmatrix}
0 & 2 & 3 & \infty & 8 & \infty & \infty & \infty \\
2 & 0 & 4 & \infty & 9 & \infty & \infty & \infty \\
3 & 4 & 0 & 7 & \infty & \infty & \infty & \infty \\
\infty & \infty & 7 & 0 & 4 & 3 & \infty & \infty \\
8 & \infty & \infty & 4 & 0 & 5 & 5 & \infty \\
\infty & 9 & \infty & 3 & 5 & 0 & 7 & 6 \\
\infty & \infty & \infty & \infty & 5 & 7 & 0 & 1 \\
\infty & \infty & \infty & \infty & \infty & 6 & 1 & 0
\end{bmatrix}
\end{equation}

Rezultat:
\begin{itemize}
\item $T = \begin{bmatrix} 1 & 2 \\ 1 & 3 \\ 3 & 4 \\ 4 & 5 \\ 4 & 6 \\ 5 & 7 \\ 7 & 8 \end{bmatrix}$
\item $V = 25$
\end{itemize}

Očekivano rješenje: $V = 25$

\subsection{Testni primjer 2 - Graf sa 4 čvora}

Ulazna težinska matrica:
\begin{equation}
W = \begin{bmatrix}
0 & 1 & 3 & 4 \\
1 & 0 & 2 & \infty \\
3 & 2 & 0 & 5 \\
4 & \infty & 5 & 0
\end{bmatrix}
\end{equation}

Rezultat:
\begin{itemize}
\item $T = \begin{bmatrix} 1 & 2 \\ 1 & 4 \\ 2 & 3 \end{bmatrix}$
\item $V = 7$
\end{itemize}

Ovo je optimalno rješenje (grane: 1-2, 1-4, 2-3 sa ukupnom težinom $1 + 4 + 2 = 7$).

\subsection{Testni primjer 3 - Kompletan graf sa 3 čvora}

Ulazna težinska matrica:
\begin{equation}
W = \begin{bmatrix}
0 & 5 & 3 \\
5 & 0 & 2 \\
3 & 2 & 0
\end{bmatrix}
\end{equation}

Rezultat:
\begin{itemize}
\item $T = \begin{bmatrix} 1 & 3 \\ 2 & 3 \end{bmatrix}$
\item $V = 5$
\end{itemize}

Optimalno rješenje (grane: 1-3, 2-3 sa ukupnom težinom $3 + 2 = 5$).

\subsection{Testni primjer 4 - Graf sa 5 čvorova}

Ulazna težinska matrica:
\begin{equation}
W = \begin{bmatrix}
0 & 2 & 3 & \infty & \infty \\
2 & 0 & 1 & 4 & \infty \\
3 & 1 & 0 & 5 & 6 \\
\infty & 4 & 5 & 0 & 7 \\
\infty & \infty & 6 & 7 & 0
\end{bmatrix}
\end{equation}

Rezultat:
\begin{itemize}
\item $T = \begin{bmatrix} 1 & 2 \\ 2 & 3 \\ 2 & 4 \\ 3 & 5 \end{bmatrix}$
\item $V = 13$
\end{itemize}

Optimalno rješenje (grane: 1-2, 2-3, 2-4, 3-5 sa ukupnom težinom $2 + 1 + 4 + 6 = 13$).

\section{Programski kod sa objašnjenjima}

U nastavku je dat kompletan programski kod sa detaljnim objašnjenjima svakog koraka:

\begin{lstlisting}[style=julia, caption={Kompletan Julia kod za minimalno povezujuće stablo}]
using JuMP
using GLPK

function min_tree(W)
    n = size(W, 1)
    
    model = Model(GLPK.Optimizer)
    
    @variable(model, x[i=1:n, j=1:n; i < j && W[i,j] < Inf], 
              Int, lower_bound=0, upper_bound=1)
    
    @objective(model, Min, 
               sum(W[i,j] * x[i,j] for i in 1:n for j in (i+1):n if W[i,j] < Inf))
    
    @constraint(model, 
                sum(x[i,j] for i in 1:n for j in (i+1):n if W[i,j] < Inf) == n - 1)
    
    svi_podskupovi = []
    for mask in 1:(2^n - 2)
        podskup = [i for i in 1:n if (mask & (1 << (i-1))) != 0]
        if length(podskup) > 0 && length(podskup) < n
            push!(svi_podskupovi, podskup)
        end
    end
    
    for S in svi_podskupovi
        komplement_S = [i for i in 1:n if !(i in S)]
        
        grane_izmedju = []
        for i in S
            for j in komplement_S
                if i < j && W[i,j] < Inf
                    push!(grane_izmedju, x[i,j])
                elseif j < i && W[j,i] < Inf
                    push!(grane_izmedju, x[j,i])
                end
            end
        end
        
        if length(grane_izmedju) > 0
            @constraint(model, sum(grane_izmedju) >= 1)
        end
    end
    
    optimize!(model)
    
    if termination_status(model) != MOI.OPTIMAL
        error("Problem nije optimalno rijesen! Status: ", termination_status(model))
    end
    
    grane = []
    for i in 1:n
        for j in (i+1):n
            if W[i,j] < Inf
                if value(x[i,j]) > 0.5
                    push!(grane, [i, j])
                end
            end
        end
    end
    
    sort!(grane, by = g -> (g[1], g[2]))
    
    if length(grane) > 0
        T = hcat([g[1] for g in grane], [g[2] for g in grane])
    else
        T = zeros(Int, 0, 2)
    end
    
    V = objective_value(model)
    
    return T, V
end
\end{lstlisting}

\subsection{Objašnjenje ključnih koraka}

\textbf{Korak 1-2:} Inicijalizacija modela - određujemo broj čvorova i kreiramo optimizacioni model koristeći GLPK solver.

\textbf{Korak 3:} Definisanje varijabli - svaka grana $(i,j)$ ima binarnu varijablu $x_{ij}$ koja označava da li je grana uključena u stablo. Varijable su cjelobrojne (Int) kako bi se osiguralo cjelobrojno rješenje.

\textbf{Korak 4:} Funkcija cilja - minimizujemo ukupnu težinu grana koje su uključene u stablo.

\textbf{Korak 5:} Ograničenje broja grana - osigurava da stablo ima tačno $n-1$ grana.

\textbf{Korak 6-7:} Generisanje ograničenja povezanosti - za svaki podskup čvorova $S$, osiguravamo da postoji najmanje jedna grana između $S$ i komplementa $S$. Ovo eliminiše cikluse i osigurava povezanost.

\textbf{Korak 8-9:} Rješavanje problema i provjera statusa.

\textbf{Korak 10-12:} Ekstrakcija rješenja - pronalazimo sve grane za koje je $x_{ij} = 1$, sortiramo ih i formiramo matricu $T$.

\textbf{Korak 13:} Izračunavanje ukupne težine stabla iz optimalne vrijednosti funkcije cilja.

\end{document}

